The quality assurance and dosimetric verification of the High Dose Rate (HDR) Brachytherapy unit was successfully completed using a well-type ionization chamber. The experimental procedure primarily aimed to determine the Reference Air Kerma Rate (RAKR) of the $^{192}$Ir source, pinpoint the precise maximum response position (the "sweet spot") within the chamber, and systematically evaluate the electromechanical accuracy of the afterloader's treatment timer.

By meticulously stepping the source through the well-chamber, the sweet spot was experimentally identified at a dwell position of 127.4 cm (Channel 10), generating a stable peak ionization current of 51.56 nA. Measurements taken at this optimal position—after applying rigorous correction factors for ambient temperature, pressure, ion recombination, and polarity—yielded an absolute RAKR of 24.1456 mGy h$^{-1}$ at 1 meter. Using the specific gamma-ray constant for $^{192}$Ir, this RAKR translates to an apparent source activity of 5.931 Ci. This result is crucial for verifying that the current clinical output matches the manufacturer’s decay-corrected specifications within strictly defined clinical tolerances (typically $\pm 3\%$).

In assessing the temporal fidelity of the HDR delivery system, the calculated intrinsic timer error was established at 0.8431 seconds. This minor discrepancy primarily accounts for the transit time required for the source to travel from the protective safe to the active dwell position and back. Furthermore, an extensive linearity check across multiple set dwell times spanning from 60 seconds to 300 seconds demonstrated an exceptionally high degree of proportionality. The calculated percentage time linearity error was a negligible 0.02\%, coupled with an end error (y-intercept) of 0.2650 seconds. 

In summary, the dosimetric calculations and temporal error analyses collectively confirm that both the physical strength of the $^{192}$Ir source and the mechanical performance of the HDR afterloader's timing mechanism fall well within internationally accepted safety and quality assurance protocols. Consequently, the unit is deemed mechanically robust, dosimetrically accurate, and fundamentally safe for high-precision clinical brachytherapy treatments.

% ==============================================================================
% SUGGESTED FIGURES TO ATTACH (Find these online or in textbooks):
%
% 1. [In Theory Section]: A schematic cross-section of an HDR Ir-192 source capsule 
%    (showing the active core, titanium encapsulation, and cable weld).
% 
% 2. [In Theory Section]: A diagram illustrating the definitions of the TG-43 
%    parameters (showing distance 'r' and polar angle 'theta' relative to the source).
%
% 3. [In Apparatus Section]: A photograph or schematic of a typical Well-Type 
%    Ionization Chamber showing the inner well and the triaxial cable connection.
%
% 4. [In Theory Section]: A visual representation of the HDR Afterloader System, 
%    showing the tungsten safe, indexer, transfer tubes, and patient applicator.
% ==============================================================================